\documentclass[12pt,a4paper]{article}
\usepackage[utf8]{inputenc}
\usepackage[english]{babel}
\usepackage{graphicx}
\usepackage{array}
\usepackage{geometry}
\usepackage{float}

\geometry{
  a4paper,
  total={170mm,257mm},
  left=20mm,
  top=20mm,
}

\newcolumntype{P}[1]{>{\hspace{0pt}}p{#1}}

\begin{document}
\begin{titlepage}
  \centering
      {\scshape\LARGE Universidad Nacional Autónoma de México \par}
      \vspace{1cm}
	     {\scshape\Large Facultad de Ciencias\par}
	     \vspace{1.2cm}
             \begin{center}
	       \includegraphics[scale=.2]{./imagenes/UNAM.eps}
	     \end{center}
	     \vspace{.5 cm}

	     {\huge\bfseries Proyecto 1: \par}
	     {\huge\bfseries Travelling Salesman Problem\par}
	     \vspace{0.5cm}

	     {\Large\itshape Emiliano Galeana Araujo\\314032324\par}
	     \vfill
	     \vspace{0.5cm}

             Análisis de Software Malicioso\par
	     \textsc{Dr. Canek Peláez Valdés}\\
             \textsc{Víctor Zamora Gutiérrez}\\
	     \vspace{0.1cm}
	            {\large \today \par}
\end{titlepage}

NOTA: se corre desde la carpeta \textbf{main} con \texttt{go run main.go}. Y solo
muestra todas las listas que hago y los valores de la máxima distancia,
normalizador y la función de costo.\\
NOTA2: Esto solo es provicional, luego lo pongo chido.
\end{document}
